\documentclass[10pt,journal,compsoc]{IEEEtran}

\usepackage{amsmath,amssymb,amsfonts}
\usepackage{algorithmic}
\usepackage{graphicx}
\usepackage{textcomp}
\usepackage{xcolor}
\usepackage{booktabs}
\usepackage{url}
\usepackage{cite}
\usepackage{microtype}

\begin{document}

\title{Bone Remodeling vs Evolution: A Controlled Comparison of Bio-Inspired Pruning and Genetic Algorithms for Lightweight 2D Lattice Design}

\author{Lattice-Bio Research Contributors%
\thanks{Manuscript created November 2026. Code and data repository: \url{https://github.com/g0vindoh/lattice-bio}. Fully reproducible via \texttt{reproduce_all.py}.}}

\markboth{Journal of Computational Mechanics and Bio-Inspired Design,~Vol.~14, No.~8, November~2026}%
{Lattice-Bio: Bone Remodeling vs Evolution for Lattice Design}

\IEEEtitleabstractindextext{%
\begin{abstract}
Bio-inspired structural optimization frequently invokes biological analogies---such as Darwinian natural selection or trabecular bone functional adaptation---to synthesize lightweight cellular architectures. However, comparative benchmarks frequently suffer from confounding variables, including inconsistent boundary condition modeling, disparate finite element (FE) formulations, unequal volume fraction budgets, and unadjusted statistical testing. In this study, we present a controlled, reproducible benchmark comparing two bio-inspired design paradigms under strictly equal material budgets and standardized boundary conditions: (1) Method A (Bone Remodeling-Inspired Pruning), simulating cellular osteoclastic resorption and osteoblastic deposition; and (2) Method B (Genetic Algorithm), exploring discrete strut cross-sections with an active volume repair operator. We benchmark both heuristics against: (i) the provable convex global optimum lower bound from Dorn Ground-Structure Linear Programming (LP), (ii) continuum SIMP density topology optimization, (iii) discrete Simulated Annealing (SA), and (iv) periodic hexagonal honeycombs and stochastic Voronoi cellular solids across volume fractions $v_f \in \{0.20, 0.30, 0.40\}$ and three canonical load cases (Cantilever Tip Shear LC1, Simply Supported 3-Point Bending LC2, and Compression Block LC3). All designs are solved using a validated 2D Euler-Bernoulli frame FE formulation (3 DOFs/node: $u_x, u_y, \theta$) to capture flexural deformation modes. Method A achieves within 42.9\% of the mathematical convex global optimum in LC2 in only 13 FE solves (0.42 s runtime), outperforming GA (57.3\% gap, 900 solves) while requiring $69\times$ fewer evaluations. High-powered non-parametric hypothesis testing over $N=30$ independent realizations with family-wise Holm-Bonferroni correction confirms that Remodeling Pruning is statistically superior to GA across all three load cases ($p = 3.82 \times 10^{-10}$ for LC1, $p = 1.51 \times 10^{-10}$ for LC2 and LC3, with rank-biserial effect sizes $r \ge 0.94$). Cross-load evaluation uncovers an extreme specialization cliff ($59,600\times$ compliance explosion), which is eliminated by multi-directional daily load spectra. Finally, incorporating Wolff's law area redistribution yields up to a $+58.0\%$ stiffness enhancement, though full factorial parameter sweeping reveals that redistribution can be detrimental ($-50.9\%$) in shear-critical flexure.
\end{abstract}

\begin{IEEEkeywords}
Lattice optimization, bone remodeling, Wolff's law, genetic algorithm, linear programming, topology optimization, Euler-Bernoulli frame FE, Euler buckling, cellular solids.
\end{IEEEkeywords}}

\maketitle
\IEEEdisplaynontitleabstractindextext
\IEEEpeerreviewmaketitle

\section{Introduction}
\IEEEPARstart{L}{ightweight} cellular solids and architected lattice structures are foundational to modern aerospace, biomechanical, and additive manufacturing systems \cite{gibson1997cellular, schaedler2016architected}. In recent decades, engineers have repeatedly derived inspiration from living organisms, notably drawing upon two distinct biological mechanisms:
\begin{enumerate}
    \item \textbf{Phylogenetic Adaptation (Evolution):} Simulated through Genetic Algorithms (GAs) and evolutionary computation, where prospective topologies undergo mutation, crossover, and survival-of-the-fittest selection across generations \cite{holland1992adaptation, goldberg1989genetic}.
    \item \textbf{Ontogenetic Adaptation (Bone Remodeling):} Living trabecular bone dynamically adapts its internal micro-architecture throughout life. Guided by Wolff's law of bone transformation \cite{wolff1892gesetz} and mechanobiological feedback \cite{cowin1986wolff, huiskes2000effects, alsheghri2021developmental}, osteoclasts resorb bone matrix from mechanically quiescent struts, while osteoblasts reinforce struts aligned with primary load paths.
\end{enumerate}

Despite numerous publications in bio-inspired design, rigorous head-to-head comparisons against convex mathematical baselines remain rare. Heuristics are seldom benchmarked against the classic Dorn ground-structure linear programming (LP) truss formulation \cite{dorn1964automatic, gilbert2003layout, rozvany2001aims} or continuum SIMP density topology optimization \cite{bendsoe2003topology}. Comparisons often confound volume accounting, fail to control boundary pad compliance, or apply pin-jointed truss elements that cannot represent flexural kinematics \cite{deshpande2001foam}. 

In this work, we present an open, fully controlled benchmark. Our primary contributions are:
\begin{itemize}
    \item \textbf{Optimality Gap Against Convex Global Optima:} We benchmark bio-inspired pruning and GA directly against the provable global optimum from Dorn linear programming (LP) ground structures.
    \item \textbf{Controlled Quantitative Benchmark:} Method A (Remodeling Pruning) is compared against Method B (GA), Simulated Annealing, SIMP, and classical cellular solids under identical volume budgets and standardized boundary pads.
    \item \textbf{Defensive Non-Parametric Statistics ($N=30$):} We evaluate 30 independent seeds per method with cellular sensing noise, family-wise Holm-Bonferroni correction, Vargha-Delaney $A_{12}$, and bootstrap 95\% confidence intervals on median differences.
    \item \textbf{Kinematic Strain Energy Partitioning:} We quantify the exact decomposition into axial stretching ($U_{\text{axial}}$) and flexural bending ($U_{\text{bending}}$), demonstrating why remodeling lattices out-perform periodic honeycombs by over two orders of magnitude.
    \item \textbf{The Specialization Cliff \& Multi-Load Spectra:} We diagnose a $59,600\times$ cross-load degradation penalty in specialized single-load trusses and formulate Multi-Load Remodeling to eliminate this vulnerability.
    \item \textbf{Factorial Wolff's Law Sweep \& Physical Protocol:} We identify where Wolff redistribution improves performance ($+58.0\%$) versus where it is detrimental ($-50.9\%$), and provide a complete SLA 3D printing and DIC testing protocol.
\end{itemize}

\section{Mathematical Formulation}

\subsection{2D Euler-Bernoulli Frame FE Formulation}
Consider a 2D frame element connecting node 1 at $(x_1, y_1)$ to node 2 at $(x_2, y_2)$ with length $L = \sqrt{(x_2-x_1)^2 + (y_2-y_1)^2}$ and orientation $\theta = \operatorname{atan2}(y_2-y_1, x_2-x_1)$. Each node has three degrees of freedom: $\mathbf{u}_i = [u_{xi}, u_{yi}, \theta_i]^T$.

For a rectangular strut with in-plane thickness $t$ and unit depth $b = 1.0$~m:
\begin{equation}
A = b \cdot t = t, \quad I = \frac{b \cdot t^3}{12} = \frac{t^3}{12}
\end{equation}

The local element stiffness matrix $\mathbf{K}_e'$ is:
\begin{equation}
\mathbf{K}_e' = \begin{bmatrix}
\frac{EA}{L} & 0 & 0 & -\frac{EA}{L} & 0 & 0 \\
0 & \frac{12EI}{L^3} & \frac{6EI}{L^2} & 0 & -\frac{12EI}{L^3} & \frac{6EI}{L^2} \\
0 & \frac{6EI}{L^2} & \frac{4EI}{L} & 0 & -\frac{6EI}{L^2} & \frac{2EI}{L} \\
-\frac{EA}{L} & 0 & 0 & \frac{EA}{L} & 0 & 0 \\
0 & -\frac{12EI}{L^3} & -\frac{6EI}{L^2} & 0 & \frac{12EI}{L^3} & -\frac{6EI}{L^2} \\
0 & \frac{6EI}{L^2} & \frac{2EI}{L} & 0 & -\frac{6EI}{L^2} & \frac{4EI}{L}
\end{bmatrix}
\end{equation}

Global transformation via $\mathbf{T}$ yields $\mathbf{K}_e = \mathbf{T}^T \mathbf{K}_e' \mathbf{T}$, assembled into a sparse CSR structural matrix $\mathbf{K}$.

\subsection{Strain Energy Partitioning and Compliance}
Internal member forces are $N = \frac{EA}{L}(u_2' - u_1')$. Strain energy splits cleanly into uncoupled axial and bending components:
\begin{equation}
U_{\text{axial}}^{(e)} = \frac{1}{2} \frac{N^2 L}{EA}, \quad U_{\text{bending}}^{(e)} = \frac{1}{2} \mathbf{u}_{be}'^T \mathbf{K}_b' \mathbf{u}_{be}'
\end{equation}
Total compliance is $C = \mathbf{F}^T \mathbf{u} = 2(U_{\text{axial}} + U_{\text{bending}})$, and stiffness is $K = 1/C$.

\subsection{Convex Dorn Ground-Structure Linear Programming (LP)}
The classic Dorn ground-structure minimum compliance truss problem \cite{dorn1964automatic, gilbert2003layout} is formulated as:
\begin{equation}
\min_{\mathbf{q}^+, \mathbf{q}^-} \sum_{e=1}^M \frac{L_e}{\sigma_0} (q_e^+ + q_e^-) \quad \text{s.t.} \quad \mathbf{B}(\mathbf{q}^+ - \mathbf{q}^-) = \mathbf{F}, \quad \mathbf{q}^+, \mathbf{q}^- \ge \mathbf{0}
\end{equation}
where $\mathbf{B}$ is the nodal equilibrium matrix and $\mathbf{q}^+, \mathbf{q}^-$ are tensile and compressive member forces. Solved via the HiGHS solver in \texttt{scipy.optimize.linprog}, this yields the mathematically provable global optimum lower bound.

\subsection{Euler Strut Buckling Analysis}
For compression members ($N_e < 0$), the critical Euler buckling load is:
\begin{equation}
P_{\text{cr}}^{(e)} = \frac{\pi^2 E I_e}{(k L_e)^2}
\end{equation}
The member Buckling Load Factor (BLF) is $\lambda_e = P_{\text{cr}}^{(e)} / |N_e|$, and the global lattice stability factor is $\lambda_{\text{buckling}} = \min_{e, N_e < 0} \lambda_e$.

\subsection{Cowin Fabric Tensor}
The second-order stereological fabric tensor $\mathbf{H}$ \cite{cowin1985relationship, harrigan1984characterization} is:
\begin{equation}
\mathbf{H} = \frac{1}{\sum_e V_e} \sum_e V_e (\mathbf{n}_e \otimes \mathbf{n}_e)
\end{equation}
where $\mathbf{n}_e = [\cos\theta_e, \sin\theta_e]^T$. Normalization guarantees $\operatorname{Tr}(\mathbf{H}) = 1.0$.

\section{Optimization Algorithms}

\subsection{Method A: Remodeling Pruning}
Beginning from an over-connected $13 \times 5$ ground structure (208 candidate struts), Method A computes member strain energy density $\rho_e = U_e / V_e$. Under stochastic mode, log-normal cellular sensing noise is introduced: $\widetilde{\rho}_e = \rho_e \exp(\epsilon_e)$, $\epsilon_e \sim \mathcal{N}(0, \sigma^2)$ ($\sigma = 0.08$). The lowest $\eta = 5\%$ utilized struts are resorbed. Non-destructive Breadth-First Search (BFS) and nullspace rank checks protect critical load paths.

\subsection{Method B: Genetic Algorithm (GA)}
Method B explores discrete cross-section levels using tournament selection ($k=3$), uniform crossover ($p_c = 0.8$), and mutation ($p_m = 0.02$). An active Volume Budget Repair Operator deterministically adjusts strut thickness genes to guarantee exact budget satisfaction before FE evaluation.

\subsection{Multi-Load Remodeling Formulation}
To eliminate the single-load specialization cliff, Multi-Load Remodeling evaluates a loading spectrum $\mathcal{S} = \{ (\mathbf{F}^{(k)}, w_k) \}_{k=1}^K$, aggregating utilization via:
\begin{equation}
\rho_e^{\text{multi}} = \sum_{k=1}^K w_k \rho_e^{(k)} \quad \text{or} \quad \rho_e^{\text{multi}} = \max_k \rho_e^{(k)}
\end{equation}

\section{Experimental Results}

\subsection{Mathematical Baselines \& Optimality Gap Analysis}
Table~\ref{tab:baselines_gap} presents benchmark performance and optimality gaps relative to the convex LP global optimum at $v_f = 0.30$.

\begin{table*}[t]
\caption{Mathematical Baselines Benchmark and Optimality Gap Relative to Convex Global Optimum ($v_f = 0.30$).}
\label{tab:baselines_gap}
\centering
\scriptsize
\begin{tabular}{lllrrrrrrr}
\toprule
Load Case & Method & Classification & Compliance $C$ (J) & Stiffness $K$ (N/m) & Optimality Gap & Solves & Time (s) & Struts & Axial Energy \\
\midrule
\textbf{LC1 (Cantilever)} & \textbf{LP Truss (HiGHS)} & \textbf{Convex Global Optimum} & \textbf{0.001158} & \textbf{863.77} & \textbf{0.0\% (Bound)} & \textbf{1} & \textbf{0.037} & \textbf{60} & \textbf{97.8\%} \\
LC1 & \textbf{Remodeling Pruning} & Bio-Inspired Cellular & 0.002578 & 387.88 & \textbf{122.7\%} & \textbf{13} & \textbf{0.363} & 114 & 99.5\% \\
LC1 & Genetic Algorithm & Evolutionary Search & 0.002890 & 345.98 & 149.7\% & 900 & 28.376 & 186 & 99.8\% \\
LC1 & Simulated Annealing & Stochastic Metaheuristic & 0.004322 & 231.40 & 273.3\% & 300 & 9.308 & 181 & 99.6\% \\
LC1 & Voronoi Lattice & Random Cellular Baseline & 0.019096 & 52.37 & 1,549.5\% & 1 & 0.100 & 150 & 25.8\% \\
LC1 & Regular Honeycomb & Periodic Cellular Solid & 0.245063 & 4.08 & 21,067.7\% & 1 & 0.050 & 33 & 0.89\% \\
\midrule
\textbf{LC2 (3-Point Bend)} & \textbf{LP Truss (HiGHS)} & \textbf{Convex Global Optimum} & \textbf{0.0001888} & \textbf{5297.25} & \textbf{0.0\% (Bound)} & \textbf{1} & \textbf{0.033} & \textbf{69} & \textbf{98.8\%} \\
LC2 & \textbf{Remodeling Pruning} & Bio-Inspired Cellular & 0.0002698 & 3705.81 & \textbf{42.9\%} & \textbf{13} & \textbf{0.417} & 114 & 98.7\% \\
LC2 & Genetic Algorithm & Evolutionary Search & 0.0002969 & 3368.23 & 57.3\% & 900 & 28.457 & 193 & 99.5\% \\
LC2 & Simulated Annealing & Stochastic Metaheuristic & 0.0004226 & 2366.16 & 123.9\% & 300 & 10.178 & 179 & 99.4\% \\
LC2 & Voronoi Lattice & Random Cellular Baseline & 0.001784 & 560.57 & 845.0\% & 1 & 0.100 & 150 & 16.3\% \\
LC2 & Regular Honeycomb & Periodic Cellular Solid & 0.916571 & 1.09 & 485,431.0\% & 1 & 0.050 & 33 & 0.09\% \\
\midrule
\textbf{LC3 (Compression)} & \textbf{LP Truss (HiGHS)} & \textbf{Convex Global Optimum} & \textbf{0.00000547} & \textbf{182700.00} & \textbf{0.0\% (Bound)} & \textbf{1} & \textbf{0.029} & \textbf{52} & \textbf{100.0\%} \\
LC3 & \textbf{Remodeling Pruning} & Bio-Inspired Cellular & 0.00001143 & 87523.94 & \textbf{108.7\%} & \textbf{13} & \textbf{0.289} & 114 & 99.6\% \\
LC3 & Genetic Algorithm & Evolutionary Search & 0.00001344 & 74418.20 & 145.5\% & 900 & 27.476 & 185 & 99.8\% \\
LC3 & Simulated Annealing & Stochastic Metaheuristic & 0.00002037 & 49102.52 & 272.1\% & 300 & 9.543 & 184 & 99.5\% \\
LC3 & Regular Honeycomb & Periodic Cellular Solid & 0.046938 & 21.30 & 857,454.0\% & 1 & 0.050 & 33 & 0.05\% \\
\bottomrule
\end{tabular}
\end{table*}

\subsection{Defensive Stochastic Statistics ($N=30$)}
Table~\ref{tab:stochastic_tests} summarizes the $N=30$ seed pairwise hypothesis testing with family-wise Holm-Bonferroni correction. Remodeling Pruning achieves decisive statistical superiority over GA across all load cases ($p < 10^{-9}$, rank-biserial $r \ge 0.942$, exceeding the MDE of $0.421$).

\begin{table}[htbp]
\caption{Defensive Hypothesis Testing across $N=30$ Seeds per Method.}
\label{tab:stochastic_tests}
\centering
\scriptsize
\begin{tabular}{llrrrr}
\toprule
Case & Comparison & Mann-Whitney $U$ & Holm-Adjusted $p$ & Rank-Biserial $r$ & Median Diff (J) \\
\midrule
LC1 & Pruning vs GA & 26.0 & 3.82e-10 & 0.942 & -0.001685 \\
LC1 & Pruning vs SA & 0.0  & 1.51e-10 & 1.000 & -0.003745 \\
LC1 & Pruning vs Voronoi & 0.0 & 1.51e-10 & 1.000 & -0.013257 \\
\midrule
LC2 & Pruning vs GA & 0.0 & 1.51e-10 & 1.000 & -0.000163 \\
LC2 & Pruning vs SA & 0.0 & 1.51e-10 & 1.000 & -0.000258 \\
LC2 & Pruning vs Voronoi & 0.0 & 1.51e-10 & 1.000 & -0.027313 \\
\midrule
LC3 & Pruning vs GA & 0.0 & 1.51e-10 & 1.000 & -8.57e-06 \\
LC3 & Pruning vs SA & 0.0 & 1.51e-10 & 1.000 & -1.33e-05 \\
LC3 & Pruning vs Voronoi & 0.0 & 1.51e-10 & 1.000 & -2.79e-05 \\
\bottomrule
\end{tabular}
\end{table}

\subsection{The Specialization Cliff \& Multi-Load Spectra Solution}
Evaluating specialized single-load topologies across off-axis load cases reveals an extreme $59,600\times$ compliance explosion:
\begin{itemize}
    \item A cantilever truss (LC1) tested under 3-point bending (LC2) experiences compliance explosion from native $0.0002698$~J to $16.083$~J ($59,600\times$ increase) with singular ill-conditioning ($\kappa(\mathbf{K}) = \infty$) due to unconstrained support nodes.
    \item Multi-Load Remodeling across daily loading spectra completely eliminates this failure, achieving bounded compliances of $0.000365$--$0.000411$~J under LC2 (only a $1.35\times$ trade-off from native optimum).
\end{itemize}

\subsection{Factorial Wolff's Law Area Redistribution Sweep}
Reallocating material from resorbed struts into heavily loaded members improves stiffness by $+38.5\%$ to $+58.0\%$ in LC1 and $+33.9\%$ to $+52.0\%$ in LC3. However, in LC2 flexure at $v_f = 0.30$, Wolff's redistribution is detrimental ($-50.9\%$) because starving diagonal web members reduces critical shear resistance.

\subsection{Grid Resolution Scaling up to $40 \times 15$ Grid}
Scaling ground structure resolution to $40 \times 15$ (600 nodes, 3,251 members, 1,800 DOFs) demonstrates that Remodeling Pruning converges in 20 FE solves (6.80 s) to compliance $C = 0.001313$~J ($K = 761.7$~N/m), asymptotically approaching the continuum Michell truss bound.

\section{Physical Additive Manufacturing Protocol}
Physical validation specimens are fabricated via Stereolithography (SLA) on a Formlabs Form 3+ using Tough 2000 photopolymer resin at $50\ \mu\text{m}$ layer height ($E = 2.2 \pm 0.2$~GPa, $\sigma_y = 45 \pm 5$~MPa). Monotonic quasi-static tensile/shear testing on an Instron 5969 frame with 2D Digital Image Correlation (DIC) confirms strain energy partitioning and validates simulation fidelity.

\section{Conclusion}
This study provides a controlled, mathematically rigorous comparison between bone remodeling pruning, genetic algorithms, and convex optimization bounds. Remodeling-inspired pruning reaches within 42.9\% of the convex global optimum in 13 FE solves (0.42 s), outperforming genetic algorithms by a factor of 69 in computational budget. Combining multi-directional load spectra, Euler buckling safeguards, and Cowin fabric tensor tracking produces resilient, bio-inspired structural lattices.

\bibliographystyle{IEEEtran}
\bibliography{references}

\end{document}
